\documentclass{article}
\usepackage{geometry}
\geometry{margin=1.5cm, vmargin={0pt,1cm}}
\setlength{\topmargin}{-1cm}
\setlength{\headheight}{12.5pt}
\setlength{\paperheight}{29.7cm}
\setlength{\textheight}{25.3cm}

% useful packages
\usepackage{fancyhdr}
\usepackage{layout}
\usepackage{tikz}
\usetikzlibrary{arrows.meta}

% import common macros
% % % % % % % % % % % % % % % % % % % % % % % % % % % % % % % %
% Common Macros for LaTeX
% % % % % % % % % % % % % % % % % % % % % % % % % % % % % % % %

% Useful packages
\usepackage{amsfonts}
\usepackage{amsmath}
\usepackage{amssymb}
\usepackage{amsthm}
\usepackage{enumerate}
\usepackage{graphicx}
\usepackage{multicol}
\usepackage[ruled,linesnumbered]{algorithm2e}

% Math operators and common symbols
\newcommand{\dif}{\mathrm{d}}
\newcommand{\innerProd}[1]{\left\langle #1 \right\rangle}
\newcommand{\difFrac}[2]{\frac{\dif #1}{\dif #2}}
\newcommand{\pdfFrac}[2]{\frac{\partial #1}{\partial #2}}
\newcommand{\T}{\mathrm{T}}
\newcommand{\norm}[1]{\left\| #1 \right\|}
\newcommand{\abs}[1]{\left| #1 \right|}

% Vector and Matrix shortcuts
\newcommand{\vecTwo}[2]{
  \begin{bmatrix}
    #1 \\
    #2
\end{bmatrix}}
\newcommand{\vecThree}[3]{
  \begin{bmatrix}
    #1 \\
    #2 \\
    #3
\end{bmatrix}}
\newcommand{\vecTwoH}[2]{
  \begin{bmatrix}
    #1 & #2
  \end{bmatrix}
}
\newcommand{\vecThreeH}[3]{
  \begin{bmatrix}
    #1 & #2 & #3
  \end{bmatrix}
}

% Common fields
\newcommand{\R}{\mathbb{R}}
\newcommand{\C}{\mathbb{C}}
\newcommand{\Z}{\mathbb{Z}}
\newcommand{\N}{\mathbb{N}}

% Solutions and proofs
\newenvironment{solution}{
\begin{proof}[解]}{
\end{proof}}

% Utility to get environment variables (requires lualatex)
\newcommand{\getEnv}[2]{\directlua{tex.print(os.getenv("#1") or "#2")}}


\providecommand{\Chat}{\overline{\C}}
\providecommand{\Mob}{\mathfrak{M}}
\providecommand{\GL}{\operatorname{GL}}
\providecommand{\SL}{\operatorname{SL}}
\providecommand{\Aut}{\operatorname{Aut}}

% Name and uID
\newcommand{\name}{\getEnv{NAME}{NAME}}
\newcommand{\uid}{\getEnv{UID}{UID}}
\newcommand{\course}{Complex Analysis}
\newcommand{\hwNum}{12}

\begin{document}

\pagestyle{fancy}
\fancyhead{}
\lhead{\zhstr{\name} (\uid)}
\chead{\course\ \#\hwNum}
\rhead{\today}

\section{Exercise 4.134}

\textbf{Lemma 4.133.}
A set $\mathcal D\subseteq \Chat$ is open if and only if the set
\[
  \{z\in\Chat:z^{-1}\in \mathcal D\}
\]
is open in $\Chat$.

\begin{solution}
  Let
  \[
    I:\Chat\to\Chat,\qquad I(z)=z^{-1},
  \]
  where $0^{-1}=\infty$ and $\infty^{-1}=0$. Then $I\circ I=\operatorname{id}$.
  It suffices to show that if $\mathcal D$ is open in $\Chat$, then
  $I^{-1}(\mathcal D)$ is open in $\Chat$.

  Put
  \[
    \widehat{\mathcal D}:=I^{-1}(\mathcal D)
    =\{z\in\Chat:z^{-1}\in\mathcal D\}.
  \]
  We verify openness by Definition~4.132.

  First consider finite points. Let $z_0\in \widehat{\mathcal D}\cap\C$.

  If $z_0\ne0$, then $z_0^{-1}\in \mathcal D\cap\C$. Since
  $\mathcal D\cap\C$ is open in $\C$ and $z\mapsto z^{-1}$ is continuous on
  $\C^\ast$, there is a small disk around $z_0$ contained in
  $\widehat{\mathcal D}\cap\C$.

  If $z_0=0$, then $\infty=0^{-1}\in\mathcal D$. Since $\mathcal D$ is open in
  $\Chat$, Definition~4.132 gives some $R>0$ such that
  \[
    \{w\in\C:|w|>R\}\subseteq \mathcal D.
  \]
  Thus whenever $0<|z|<1/R$, we have $|z^{-1}|>R$, so $z^{-1}\in\mathcal D$.
  Also $0^{-1}=\infty\in\mathcal D$. Hence a disk around $0$ is contained in
  $\widehat{\mathcal D}\cap\C$.

  Therefore $\widehat{\mathcal D}\cap\C$ is open in $\C$.

  It remains to check the neighborhood condition at $\infty$. If
  $\infty\in\widehat{\mathcal D}$, then $0=\infty^{-1}\in\mathcal D$. Since
  $\mathcal D\cap\C$ is open in $\C$, there exists $r>0$ such that
  \[
    U_r(0)\subseteq \mathcal D.
  \]
  Then for $|z|>1/r$, we have $|z^{-1}|<r$, hence $z^{-1}\in\mathcal D$.
  Therefore
  \[
    \{z\in\C:|z|>1/r\}\subseteq \widehat{\mathcal D}.
  \]
  So $\widehat{\mathcal D}$ is open in $\Chat$.

  Conversely, if $\widehat{\mathcal D}$ is open, then
  \[
    \mathcal D=I^{-1}(\widehat{\mathcal D})
  \]
  is open by the same argument, because $I^{-1}=I$. Hence the equivalence is
  proved.
\end{solution}

\section{Exercise 4.153}

\textbf{Theorem 4.152.}
A function $f:\Chat\to\Chat$ is meromorphic if and only if it is rational,
i.e. of the form
\[
  f(z)=\frac{P(z)}{Q(z)},
\]
where $P$ and $Q$ are complex polynomials.

\textbf{Problem.}
Give an alternative proof of the necessity part of Theorem~4.152 using
Laurent expansions and Liouville's theorem~3.94.

\begin{solution}
  Assume $f:\Chat\to\Chat$ is meromorphic. We prove that $f$ is rational.

  Since $f$ is meromorphic on $\Chat$, its pole set is discrete in $\Chat$.
  The Riemann sphere $\Chat$ is compact, so a discrete subset of $\Chat$ must be
  finite. Hence $f$ has only finitely many finite poles, say
  \[
    z_1,z_2,\dots,z_m\in\C,
  \]
  with pole orders $k_1,k_2,\dots,k_m$.

  Define
  \[
    Q(z):=\prod_{j=1}^m (z-z_j)^{k_j}.
  \]
  Then
  \[
    F(z):=Q(z)f(z)
  \]
  has removable singularities at all finite poles of $f$. After removing those
  singularities, $F$ is entire on $\C$.

  It remains to show that $F$ is a polynomial. Since $f$ is meromorphic at
  $\infty$, the function
  \[
    \widehat f(\zeta):=f(\zeta^{-1})
  \]
  has a non-essential singularity at $\zeta=0$. Also $Q(\zeta^{-1})$ has at
  worst a pole at $\zeta=0$. Therefore
  \[
    \widehat F(\zeta):=F(\zeta^{-1})=Q(\zeta^{-1})f(\zeta^{-1})
  \]
  has at worst a pole at $\zeta=0$. Thus the Laurent expansion of
  $\widehat F$ at $0$ has only finitely many negative powers.

  Equivalently, the Laurent expansion of $F$ at $\infty$ has only finitely many
  positive powers. Hence there exists a polynomial $P$ such that
  \[
    F(z)-P(z)\to 0
    \qquad (z\to\infty).
  \]
  The function $F-P$ is entire and bounded outside a sufficiently large disk;
  it is also bounded on the remaining compact disk. Thus $F-P$ is a bounded
  entire function. By Liouville's theorem~3.94, $F-P$ is constant. Since
  $F(z)-P(z)\to0$ as $z\to\infty$, this constant is $0$. Hence
  \[
    F=P.
  \]
  Therefore
  \[
    f(z)=\frac{P(z)}{Q(z)}
  \]
  is rational. This proves the necessity part of Theorem~4.152.
\end{solution}

\section{Exercise 4.155*}

Let $f$ be a monic polynomial of degree $d$, and
\[
  M(f):=\max_{z\in\partial E}|f(z)|.
\]
Show that

\begin{enumerate}
  \item[(a)] $\forall |z|>1,\ |f(z)|\le |z|^dM(f)$.
  \item[(b)] $M(f)\ge1$; the equality holds if and only if $f(z)=z^d$.
\end{enumerate}

\begin{solution}
  Since $f$ is monic of degree $d$, write
  \[
    f(z)=z^d+a_{d-1}z^{d-1}+\cdots+a_1z+a_0.
  \]
  Define
  \[
    g(w):=w^d f(w^{-1}).
  \]
  Then
  \[
    g(w)=1+a_{d-1}w+\cdots+a_1w^{d-1}+a_0w^d,
  \]
  so $g$ is analytic in a neighborhood of $\overline E$ and $g(0)=1$.

  For $|z|>1$, put $w=z^{-1}$. Then $|w|<1$, and
  \[
    |f(z)|=|z|^d |z^{-d}f(z)|
    =|z|^d |g(z^{-1})|.
  \]
  By the maximum modulus principle applied to $g$ on $E$,
  \[
    |g(z^{-1})|
    \le \max_{|w|=1}|g(w)|.
  \]
  For $|w|=1$,
  \[
    |g(w)|=|w^d f(w^{-1})|=|f(w^{-1})|.
  \]
  Since $w^{-1}\in\partial E$, we get
  \[
    \max_{|w|=1}|g(w)|=M(f).
  \]
  Therefore
  \[
    |f(z)|\le |z|^dM(f),
  \]
  which proves (a).

  For (b), since $g(0)=1$, the maximum modulus principle gives
  \[
    1=|g(0)|\le \max_{|w|=1}|g(w)|=M(f).
  \]
  Thus $M(f)\ge1$.

  If $M(f)=1$, then $|g(0)|=1$ equals the maximum modulus of $g$ on $E$.
  Hence $g$ attains its maximum modulus at an interior point. By the maximum
  modulus theorem, $g$ is constant. Since $g(0)=1$, we have $g\equiv1$.
  Therefore
  \[
    w^d f(w^{-1})=1,
  \]
  so $f(z)=z^d$.

  Conversely, if $f(z)=z^d$, then clearly
  \[
    M(f)=\max_{|z|=1}|z^d|=1.
  \]
  This completes the proof.
\end{solution}

\section{Exercise 4.158}

Show that the singularity of $f$ is a pole in the proof of Theorem~4.156 by
applying the Casorati-Weierstrass theorem~4.108 and the open mapping theorem
4.41 to $f(1/z)$.

\begin{solution}
  In the proof of Theorem~4.156, $f:\C\to\C$ is assumed to be globally
  conformal. Thus $f$ is entire, bijective, and nonconstant.

  Consider
  \[
    g(z):=f(1/z),
  \]
  which is analytic in a punctured disk $0<|z|<r$. We prove that the isolated
  singularity of $g$ at $0$ is a pole.

  First, $g$ cannot have an essential singularity at $0$. Suppose it did. By
  the Casorati-Weierstrass theorem~4.108, for every $r>0$,
  \[
    g(U_r^\bullet(0))
  \]
  is dense in $\C$. Equivalently, for every $R>0$,
  \[
    f(\{z\in\C:|z|>R\})
  \]
  is dense in $\C$.

  On the other hand, since $f$ is nonconstant and analytic, the open mapping
  theorem~4.41 implies that
  \[
    f(U_R(0))
  \]
  is a nonempty open subset of $\C$. A dense subset of $\C$ must intersect this
  open set. Hence
  \[
    f(\{z:|z|>R\})\cap f(U_R(0))\ne\varnothing.
  \]
  This contradicts the injectivity of $f$, because the two subsets
  $\{z:|z|>R\}$ and $U_R(0)$ are disjoint. Therefore $g$ is not essential at
  $0$.

  Next, $g$ cannot have a removable singularity at $0$. Suppose $g$ extends
  analytically to $0$, and denote the extension by $G$. The extension $G$ is
  nonconstant; otherwise $f$ would be constant outside a disk, hence constant
  on $\C$ by the identity theorem. Thus the open mapping theorem implies that
  $G(U_r(0))$ is open and contains $G(0)$. Let $a=G(0)$. Since $f$ is
  surjective, choose $z_0\in\C$ such that
  \[
    f(z_0)=a.
  \]
  Choose $r>0$ sufficiently small and then choose $R$ such that
  \[
    |z_0|<R<1/r.
  \]
  The set $f(U_R(0))$ is open and contains $a$. Also $G(U_r(0))$ is open and
  contains $a$. Thus their intersection contains some point $b\ne a$.

  Since $b\in G(U_r(0))$ and $b\ne a=G(0)$, there exists some
  $\zeta\in U_r^\bullet(0)$ such that
  \[
    b=G(\zeta)=f(1/\zeta),
  \]
  where $|1/\zeta|>1/r>R$. Since $b\in f(U_R(0))$, there also exists
  $w\in U_R(0)$ such that
  \[
    b=f(w).
  \]
  Thus $f(w)=f(1/\zeta)$ with $w\in U_R(0)$ and $|1/\zeta|>R$, contradicting
  the injectivity of $f$.

  Therefore the singularity of $g(z)=f(1/z)$ at $0$ is neither removable nor
  essential. By the classification of isolated singularities, it must be a
  pole. Hence the singularity of $f$ at $\infty$ is a pole.
\end{solution}

\section{Exercise 4.159}

Find all entire functions $f$ satisfying
\[
  f(f(z))=z
  \qquad \text{for all }z\in\C.
\]

\begin{solution}
  The equation
  \[
    f(f(z))=z
  \]
  implies that $f$ is injective. Indeed, if $f(z_1)=f(z_2)$, then applying $f$
  to both sides gives
  \[
    z_1=f(f(z_1))=f(f(z_2))=z_2.
  \]
  It also implies that $f$ is surjective: for any $w\in\C$,
  \[
    w=f(f(w)).
  \]
  Hence every $w$ lies in the image of $f$.

  Therefore $f$ is an entire bijection. By Theorem~4.156, $f$ must be affine:
  \[
    f(z)=az+b,
    \qquad a\in\C^\ast,\ b\in\C.
  \]
  Now compute
  \[
    f(f(z))=a(az+b)+b=a^2z+(a+1)b.
  \]
  Since $f(f(z))=z$ for all $z$, we need
  \[
    a^2=1,
    \qquad
    (a+1)b=0.
  \]
  Thus $a=1$ or $a=-1$.

  If $a=1$, then $2b=0$, so $b=0$, giving
  \[
    f(z)=z.
  \]
  If $a=-1$, then $(a+1)b=0$ is automatic, so $b$ is arbitrary, giving
  \[
    f(z)=b-z,
    \qquad b\in\C.
  \]
  Therefore all solutions are
  \[
    \boxed{f(z)=z}
    \qquad\text{and}\qquad
    \boxed{f(z)=b-z,\ b\in\C}.
  \]
\end{solution}

\section{Exercise 4.165}

Explain how to normalize any nonsingular Mobius transformation $M$.

\begin{solution}
  Let
  \[
    M(z)=\frac{az+b}{cz+d}
  \]
  be a nonsingular Mobius transformation. By Definition~4.162,
  \[
    ad-bc\ne0.
  \]
  Multiplying all coefficients by the same nonzero complex number does not
  change the Mobius transformation:
  \[
    \frac{az+b}{cz+d}
    =
    \frac{\lambda az+\lambda b}{\lambda cz+\lambda d},
    \qquad \lambda\in\C^\ast.
  \]
  We want the new coefficients to satisfy Definition~4.164, namely
  \[
    (\lambda a)(\lambda d)-(\lambda b)(\lambda c)=1.
  \]
  This condition is
  \[
    \lambda^2(ad-bc)=1.
  \]
  Since $ad-bc\ne0$, we can choose
  \[
    \lambda=(ad-bc)^{-1/2},
  \]
  where either square root may be chosen. Then
  \[
    (\lambda a)(\lambda d)-(\lambda b)(\lambda c)=1.
  \]
  Thus every nonsingular Mobius transformation can be normalized by multiplying
  all four coefficients by a common factor $\lambda$ satisfying
  \[
    \lambda^2(ad-bc)=1.
  \]
  The two possible choices differ by a sign.
\end{solution}

\section{Exercise 4.173}

Prove

\begin{enumerate}
  \item[(i)] $[f^{-1}(z)]=[f]^{-1}$;
  \item[(ii)] $[f_2][f_1]=[f_2\circ f_1]$
\end{enumerate}
for two nonsingular Mobius transformations $f_1(z),f_2(z)$.

\begin{solution}
  Let
  \[
    f(z)=\frac{az+b}{cz+d},
    \qquad
    [f]=
    \begin{bmatrix}
      a & b \\
      c & d
    \end{bmatrix},
  \]
  where $ad-bc\ne0$.

  For (i), the inverse matrix is
  \[
    [f]^{-1}
    =
    \frac1{ad-bc}
    \begin{bmatrix}
      d & -b \\
      -c & a
    \end{bmatrix}.
  \]
  Since multiplying all coefficients by the same nonzero scalar does not
  change the corresponding Mobius transformation, this matrix represents
  \[
    z\mapsto
    \frac{dz-b}{-cz+a}.
  \]
  A direct calculation shows that this is the inverse of
  \[
    z\mapsto \frac{az+b}{cz+d}.
  \]
  Indeed, if $w=f(z)$, then
  \[
    w=\frac{az+b}{cz+d}
    \implies
    w(cz+d)=az+b
    \implies
    (cw-a)z=b-dw,
  \]
  so
  \[
    z=\frac{dw-b}{-cw+a}.
  \]
  Therefore
  \[
    [f^{-1}]=[f]^{-1}.
  \]

  For (ii), write
  \[
    f_1(z)=\frac{a_1z+b_1}{c_1z+d_1},
    \qquad
    f_2(z)=\frac{a_2z+b_2}{c_2z+d_2}.
  \]
  Then
  \[
    [f_1]=
    \begin{bmatrix}
      a_1 & b_1 \\
      c_1 & d_1
    \end{bmatrix},
    \qquad
    [f_2]=
    \begin{bmatrix}
      a_2 & b_2 \\
      c_2 & d_2
    \end{bmatrix}.
  \]
  Compute
  \[
    f_2(f_1(z))
    =
    \frac{
      a_2\frac{a_1z+b_1}{c_1z+d_1}+b_2
    }{
      c_2\frac{a_1z+b_1}{c_1z+d_1}+d_2
    }.
  \]
  Multiplying numerator and denominator by $c_1z+d_1$, we get
  \[
    f_2(f_1(z))
    =
    \frac{
      (a_2a_1+b_2c_1)z+(a_2b_1+b_2d_1)
    }{
      (c_2a_1+d_2c_1)z+(c_2b_1+d_2d_1)
    }.
  \]
  Thus
  \[
    [f_2\circ f_1]
    =
    \begin{bmatrix}
      a_2a_1+b_2c_1 & a_2b_1+b_2d_1 \\
      c_2a_1+d_2c_1 & c_2b_1+d_2d_1
    \end{bmatrix}
    =
    [f_2][f_1].
  \]
\end{solution}

\section{Exercise 4.176}

Prove Lemma~4.167 using Lemma~4.175.

\textbf{Hint.} Use
\[
  \det([g][f])=\det[g]\det[f].
\]

\begin{solution}
  Lemma~4.167 states that if $f,g$ are two Mobius transformations, then so is
  $g\circ f$.

  By Lemma~4.175, the map
  \[
    \psi:\GL_2(\C)\to\Mob
  \]
  defined by
  \[
    \psi([f])=f
  \]
  is a group homomorphism from $\GL_2(\C)$ onto the Mobius group $\Mob$.

  Let $f,g\in\Mob$. Then there exist matrices
  \[
    A=[f]\in\GL_2(\C),
    \qquad
    B=[g]\in\GL_2(\C)
  \]
  such that
  \[
    \psi(A)=f,
    \qquad
    \psi(B)=g.
  \]
  Since $A,B\in\GL_2(\C)$, we have
  \[
    \det A\ne0,
    \qquad
    \det B\ne0.
  \]
  By the hint,
  \[
    \det(BA)=\det B\det A\ne0.
  \]
  Hence
  \[
    BA\in\GL_2(\C).
  \]
  Using the homomorphism property of $\psi$ from Lemma~4.175,
  \[
    \psi(BA)=\psi(B)\circ\psi(A)=g\circ f.
  \]
  Since $BA\in\GL_2(\C)$, $\psi(BA)$ is a nonsingular Mobius transformation.
  Therefore $g\circ f$ is a Mobius transformation. This proves Lemma~4.167.
\end{solution}

\section{Exercise 4.182}

Prove the sufficiency part of Lemma~4.181.

\textbf{Lemma 4.181.}
Two points $u,v\in S^2$ are antipodal if and only if their images
$p:=\sigma(u)$ and $q:=\sigma(v)$ under the stereographic projection are
related as
\[
  p=-\frac1{\overline q}.
\]

\begin{solution}
  We prove the sufficiency. Assume
  \[
    p=-\frac1{\overline q}.
  \]
  Equivalently, for finite nonzero $p,q$,
  \[
    q=-\frac1{\overline p}.
  \]

  Recall the inverse stereographic projection:
  \[
    \sigma^{-1}(z)
    =
    \left(
      \frac{2z}{1+|z|^2},
      \frac{|z|^2-1}{1+|z|^2}
    \right),
    \qquad z\in\C,
  \]
  and
  \[
    \sigma^{-1}(\infty)=N=(0,1).
  \]

  First suppose $p,q\in\C^\ast$. Let
  \[
    u=\sigma^{-1}(p),
    \qquad
    v=\sigma^{-1}(q).
  \]
  Then
  \[
    u=
    \left(
      \frac{2p}{1+|p|^2},
      \frac{|p|^2-1}{1+|p|^2}
    \right).
  \]
  Since $q=-1/\overline p$, we have $|q|=1/|p|$. Therefore
  \[
    \frac{2q}{1+|q|^2}
    =
    \frac{2(-1/\overline p)}{1+1/|p|^2}
    =
    -\frac{2p}{1+|p|^2},
  \]
  and
  \[
    \frac{|q|^2-1}{1+|q|^2}
    =
    \frac{1/|p|^2-1}{1+1/|p|^2}
    =
    -\frac{|p|^2-1}{1+|p|^2}.
  \]
  Hence
  \[
    v=-u.
  \]
  Thus $u$ and $v$ are antipodal.

  It remains to check the exceptional cases. If $p=\infty$, then the relation
  $p=-1/\overline q$ forces $q=0$. Then
  \[
    u=\sigma^{-1}(\infty)=N=(0,1),
    \qquad
    v=\sigma^{-1}(0)=(0,-1),
  \]
  so $v=-u$. The case $q=\infty$ is symmetric.

  Therefore in every case, the relation
  \[
    p=-\frac1{\overline q}
  \]
  implies that $u$ and $v$ are antipodal. This proves the sufficiency part of
  Lemma~4.181.
\end{solution}

\section{Exercise 4.187}

For a Mobius transformation $M(z)$ with two fixed points $p\ne q$, show that
if
\[
  M(z_0)=\infty
  \qquad\text{and}\qquad
  M(\infty)=w_0,
\]
then

\begin{enumerate}
  \item[(a)] $z_0+w_0=p+q$;
  \item[(b)] $\displaystyle M(z)=\frac{w_0z-pq}{z-z_0}$.
\end{enumerate}

\begin{solution}
  Write
  \[
    M(z)=\frac{az+b}{cz+d}.
  \]
  Since $M(z_0)=\infty$, the denominator vanishes at $z_0$, so
  \[
    cz_0+d=0.
  \]
  In particular, $c\ne0$. Also
  \[
    M(\infty)=\frac ac=w_0.
  \]
  Dividing numerator and denominator by $c$, we may write
  \[
    M(z)=\frac{w_0z+\beta}{z-z_0}
  \]
  for some $\beta\in\C$.

  Since $p$ and $q$ are fixed points of $M$, they are the two roots of
  \[
    z=M(z).
  \]
  Using the above form of $M$, this equation becomes
  \[
    z=\frac{w_0z+\beta}{z-z_0}.
  \]
  Equivalently,
  \[
    z(z-z_0)=w_0z+\beta,
  \]
  so
  \[
    z^2-(z_0+w_0)z-\beta=0.
  \]
  The two roots of this quadratic are $p$ and $q$. Therefore, by Vieta's
  formulas,
  \[
    p+q=z_0+w_0
  \]
  and
  \[
    pq=-\beta.
  \]
  This proves (a), and also gives
  \[
    \beta=-pq.
  \]
  Substituting this into the expression for $M$, we obtain
  \[
    M(z)=\frac{w_0z-pq}{z-z_0}.
  \]
  This proves (b).
\end{solution}

\section{Exercise 4.189}

Prove Lemma~4.188.

\textbf{Lemma 4.188.}
Let $z=z_1/z_2\in\Chat$. Then $z$ is a fixed point of a Mobius
transformation $M$ if and only if
\[
  [z_1,z_2]^T
\]
is an eigenvector of the coefficient matrix $[M]$.

\begin{solution}
  Let
  \[
    M(z)=\frac{az+b}{cz+d},
    \qquad
    [M]=
    \begin{bmatrix}
      a & b \\
      c & d
    \end{bmatrix}.
  \]
  Let
  \[
    v=
    \begin{bmatrix}
      z_1 \\
      z_2
    \end{bmatrix}
    \ne0
  \]
  be homogeneous coordinates for $z$, so that
  \[
    z=\frac{z_1}{z_2}.
  \]

  The Mobius transformation $M$ sends $z$ to
  \[
    M(z)
    =
    \frac{az_1+bz_2}{cz_1+dz_2}.
  \]
  In homogeneous coordinates, this means that $M(z)$ is represented by
  \[
    [M]v
    =
    \begin{bmatrix}
      a & b \\
      c & d
    \end{bmatrix}
    \begin{bmatrix}
      z_1 \\
      z_2
    \end{bmatrix}
    =
    \begin{bmatrix}
      az_1+bz_2 \\
      cz_1+dz_2
    \end{bmatrix}.
  \]

  Now $z$ is a fixed point of $M$ if and only if $M(z)=z$. In homogeneous
  coordinates, two nonzero vectors represent the same point of $\Chat$ if and
  only if they differ by multiplication by a nonzero scalar. Hence
  \[
    M(z)=z
  \]
  if and only if there exists $\lambda\in\C^\ast$ such that
  \[
    [M]v=\lambda v.
  \]
  But this is exactly the statement that $v=[z_1,z_2]^T$ is an eigenvector of
  $[M]$.

  Therefore $z=z_1/z_2$ is a fixed point of $M$ if and only if
  $[z_1,z_2]^T$ is an eigenvector of $[M]$.
\end{solution}

\end{document}
